%% ============================================================
%%  main.tex — структура документа
%%
%%  СТУДЕНТ ЗМІНЮЄ:
%%    1. thesis_info.tex  — свої дані (ОБОВ'ЯЗКОВО)
%%    2. chapters/*.md    — текст розділів
%%    3. references.bib   — список літератури
%%
%%  Цей файл чіпати тільки щоб додати/видалити розділ.
%%
%%  ПЕРЕД ЗБІРКОЮ:  bash tools/check.sh
%%  ЗБІРКА:         latexmk main.tex   (налаштування у latexmkrc)
%% ============================================================
\documentclass[14pt,a4paper,oneside]{extreport}
\usepackage{unithesis}
\addbibresource{references.bib}
%% ============================================================
%%  thesis_info.tex — ЄДИНИЙ ФАЙЛ ДЛЯ ЗАПОВНЕННЯ
%% ============================================================

%% ══════════════════════════════════════════════════════
%%  БЛОК 1 — Титульна сторінка
%% ══════════════════════════════════════════════════════

%% Повна назва роботи
\renewcommand{\ThesisTitle}{Вкажіть назву роботи}

%% Прізвище Ім'я По-батькові
\renewcommand{\ThesisAuthor}{Вкажіть прізвище ім'я по-батькові}

%% Посада, науковий ступінь, Прізвище І.П.
\renewcommand{\ThesisSupervisor}{Вкажіть наукового керівника}

%% Офіційна назва кафедри
\renewcommand{\ThesisDept}{Вкажіть кафедру}

%% Офіційна назва факультету або інституту
\renewcommand{\ThesisFaculty}{Вкажіть факультет}

%% Офіційна назва університету
\renewcommand{\ThesisUniversity}{Вкажіть університет}

%% Код та назва спеціальності: 122 "Комп'ютерні науки"
\renewcommand{\ThesisSpecialty}{Вкажіть спеціальність}

%% Місто та рік захисту
\renewcommand{\ThesisCity}{Київ}
\renewcommand{\ThesisYear}{2025}

%% ══════════════════════════════════════════════════════
%%  БЛОК 2 — Реферат та вступ (генеруються автоматично)
%% ══════════════════════════════════════════════════════

%% Ключові слова ВЕЛИКИМИ ЛІТЕРАМИ через кому (5–8 слів)
\renewcommand{\ThesisKeywords}{ВКАЖІТЬ КЛЮЧОВІ СЛОВА}

%% Мета роботи — починати з інфінітива ("розробити...", "дослідити...")
\renewcommand{\ThesisMeta}{вкажіть мету роботи}

%% Об'єкт дослідження — процес, явище або система
\renewcommand{\ThesisObject}{вкажіть об'єкт дослідження}

%% Предмет дослідження — методи, моделі, характеристики
\renewcommand{\ThesisSubject}{вкажіть предмет дослідження}

%% Методи дослідження — перерахування через кому
\renewcommand{\ThesisMethods}{вкажіть методи дослідження}

%% Наукова новизна — що вперше запропоновано або удосконалено
\renewcommand{\ThesisNovelty}{вкажіть наукову новизну}

%% Практичне значення отриманих результатів
\renewcommand{\ThesisPractical}{вкажіть практичне значення}

%% ══════════════════════════════════════════════════════
%%  БЛОК 3 — Необов'язкові поля
%%  Якщо не потрібно — залиште порожніми {}
%% ══════════════════════════════════════════════════════

%% Де доповідались результати (конференції, семінари)
\renewcommand{\ThesisAprobation}{}

%% Публікації за темою роботи
\renewcommand{\ThesisPublications}{}

%% ══════════════════════════════════════════════════════
%%  СТАТИСТИКА (рисунки/таблиці/джерела/додатки)
%%
%%  Рахується АВТОМАТИЧНО — нічого вписувати не треба.
%%  Значення з'являються після 2-ї компіляції (latexmk робить це сам).
%%  Якщо у рефераті стоять "?" — просто перекомпілюйте ще раз.
%% ══════════════════════════════════════════════════════

\begin{document}
\selectlanguage{ukrainian}

%% ── Титульна сторінка ────────────────────────────────────────
\makeThesisTitlePage

%% ── Реферат (генерується з thesis_info.tex) ──────────────────
%% Лічильники рис./табл./джерел/додатків — автоматичні.
%% Для коректних значень потрібна 2-га компіляція (latexmk робить сам).
\makeThesisAbstract

%% ── Зміст ───────────────────────────────────────────────────
\tableofcontents

%% ── Перелік скорочень ───────────────────────────────────────
\unnumberedchapter{Перелік умовних скорочень}
\begin{acronym}
  \acro{ПЗ}{програмне забезпечення}
  \acro{БД}{база даних}
  \acro{МН}{машинне навчання}
  \acro{НМ}{нейронна мережа}
  \acro{CNN}{Convolutional Neural Network --- згорткова нейронна мережа}
  \acro{ViT}{Vision Transformer}
  \acro{API}{Application Programming Interface}
\end{acronym}

%% ── Вступ ───────────────────────────────────────────────────
%% Актуальність + довільний текст → chapters/00_intro.md
%% Структурні поля (мета/об'єкт/предмет/новизна) → автоматично
\unnumberedchapter{Вступ}
\markdownInput{chapters/00_intro.md}
\makeThesisIntroFields

%% ── Розділи (текст у chapters/*.md) ─────────────────────────
\markdownInput{chapters/01_chapter.md}
\markdownInput{chapters/02_chapter.md}
\markdownInput{chapters/03_chapter.md}

%% ── Висновки ────────────────────────────────────────────────
\unnumberedchapter{Висновки}
\markdownInput{chapters/00_conclusions.md}

%% ── Список літератури ───────────────────────────────────────
\printbibliography[heading=bibintoc, title={Список використаних джерел}]

%% ── Додатки ─────────────────────────────────────────────────
\begin{appendices}
  \chapter{Програмний код}
  \markdownInput{chapters/appendix_a.md}
\end{appendices}

\end{document}